\documentclass{article}
\usepackage{amsmath}
\usepackage{amssymb}
\usepackage{graphicx}
\usepackage{hyperref}
\usepackage{geometry}
\geometry{margin=1in}

\begin{document}

\title{Problem Set 3}
\author{Tate Mason}
\date{\today}
\maketitle

\section*{Problem 1: Data Overview and Analysis}

\subsection*{Part A: Data Summary}

This section will serve to introduce the data. We were provided a dataset consisting of 71 markets, each witth 4 restaurants.
There are 1000 consumers in each market. The variables provided are as follows:

\begin{itemize}
  \item \textbf{ave\_dist}: Average distance to the restaurant from consumers (in miles)
  \item \textbf{mc}: Marginal cost of the restaurant
  \item \textbf{p}: Price of the restaurant
  \item \textbf{s}: Market share of the restaurant
  \item \textbf{x}: Yelp rating of the restaurant (out of 5)
\end{itemize}

I have compiled summary statistics for these variables in Table \ref{tab:summ_stats}.

\begin{table}[htbp]
    \centering
    \caption{Summary Statistics}
    \label{tab:summ_stats}
    \begin{tabular}{lcccccccc}
    \hline\hline
    Variable & N & Mean & SD & Min & P25 & Median & P75 & Max \\
    \hline
    ave\_dist & 284 & 0.845 & 0.899 & 0.014 & 0.343 & 0.546 & 1.159 & 6.543 \\
    mc        & 284 & 0.516 & 0.289 & 0.004 & 0.273 & 0.549 & 0.765 & 0.996 \\
    p         & 284 & 5.460 & 0.906 & 4.168 & 4.861 & 5.224 & 5.755 & 9.557 \\
    s         & 284 & 0.174 & 0.120 & 0.007 & 0.086 & 0.148 & 0.224 & 0.564 \\
    x         & 284 & 2.568 & 1.478 & 0.030 & 1.222 & 2.548 & 3.934 & 4.993 \\
    \hline\hline
    \end{tabular}
\end{table}

We can see that the average distance to restaurants is 0.845 miles, with a standard deviation of 0.899 miles. The average marginal cost is 0.516, and the average price is 5.460.
The market share has a mean of 0.174, and the average yelp rating is 2.568 out of 5.

\subsection*{Part B: Graphical Analysis}

\begin{center}
  \includegraphics[width=0.7\textwidth]{outputs/price_vs_share.pdf}
\end{center}

\begin{center}
  \includegraphics[width=0.7\textwidth]{outputs/distance_vs_share.pdf}
\end{center}

\begin{center}
  \includegraphics[width=0.7\textwidth]{outputs/rating_vs_share.pdf}
\end{center}

From the graphs, we observe that as the price increases, the market share tends to decrease, indicating a negative relationship between price and market share.
For the other two, the relationships become less clear, suggesting these are not strongly correlated factors in determining market share.

\section*{Problem 2: Estimation Exercises}

\begin{table}[htbp]
\centering
\caption{Estimation Results by Method}
\label{tab:methods_combined}
\begin{tabular}{lcccccc}
\hline\hline
\textbf{Method} & \textbf{Variable} & \textbf{Coef.} & \textbf{Std. Error} & \textbf{t} & \textbf{p-value} & \textbf{95\% CI} \\
\hline
\multicolumn{7}{l}{\textbf{OLS}} \\
& (Intercept) & -2.23562 & 0.33886 & -6.60 & $<1e{-}09$ & [-2.90264,\ -1.56859] \\
& x           & 0.024112 & 0.0386086 & 0.62 & 0.5328 & [-0.0518869,\ 0.100111] \\
& p           & 0.487698 & 0.0630017 & 7.74 & $<1e{-}12$ & [0.363683,\ 0.611713] \\
\hline
\multicolumn{7}{l}{\textbf{IV (MC)}} \\
& (Intercept) & 7.13029 & 1.78457 & 4.00 & $<1e{-}04$ & [3.61747,\ 10.6431] \\
& x           & 0.295728 & 0.0654541 & 4.52 & $<1e{-}05$ & [0.166885,\ 0.42457] \\
& p\_hat      & -1.35532 & 0.350383 & -3.87 & 0.0001 & [-2.04503,\ -0.66561] \\
\hline
\multicolumn{7}{l}{\textbf{IV (MC + Avg Distance)}} \\
& (Intercept) & 7.05279 & 1.78454 & 3.95 & $<1e{-}04$ & [3.54002,\ 10.5656] \\
& x           & 0.29348 & 0.0654674 & 4.48 & $<1e{-}04$ & [0.164611,\ 0.422349] \\
& p\_hat2     & -1.34007 & 0.350377 & -3.82 & 0.0002 & [-2.02976,\ -0.650371] \\
\hline\hline
\end{tabular}
\end{table}

\subsection*{Comparison of Estimation Techniques}

In this section, we examine the case of complete information on the part of the consumer. We estimate market share
using three methods: OLS, IV (MC), and IV (MC + Avg Distance). From the results in Table \ref{tab:methods_combined},
we compare the estimated coefficients for price ($\hat{\alpha}$) across the three methods. In OLS, we find a positive
coefficient, a result we know to be biased. When using IV methodologies, we find negative coefficients, which is 
at least a step in the right direction. However, the magnitudes of these estimates are large, which leads one
to question the validity of the instruments used, especially since marginal cost and average distance may not be
strongly correlated with price, a fact seen for distance in Problem 1.

None of the methods got close to the true value of -0.25 for the price coefficient, but the IV methods did at least get the sign correct.
Same goes for the rating coefficient, where OLS failed to find significance, while both IV methods found positive and significant estimates.
Finally, the intercepts were way off in all three methods, suggesting that none of these linear models are particularly well-suited
to capture the true underlying relationship.

\subsection*{Part A: OLS}

For this part, we estimate the market share equation using OLS:
\begin{gather*}
  s = \beta_0 + \beta_1 x + \beta_2 p + \varepsilon
\end{gather*}

This is the simplest case. To arrive at this equation, we formed a simple linear regression model
which reflects how market share is affected by Yelp rating (x) and price (p).

from the estimated coefficients in Table \ref{tab:methods_combined},
we find that $\hat{\alpha} = 0.487$, and is significant. However, this 
estimate is definitely biased due to price endogeneity, which can even
be seen by looking at the scatter plot of price vs. market share in Problem 1.
When price increases, market share decreases.


\subsection*{Part B: IV Estimation with Marginal Cost}

We now instrument for price using the marginal cost variable (mc).
The estimated equation is:
\begin{gather*}
  s = \beta_0 + \beta_1 x + \beta_2 \hat{p} + \varepsilon \\
  \shortintertext{where} \\
  \hat{p} = \gamma_0 + \gamma_1 x + \gamma_2 mc + \nu
\end{gather*}

Here, we regress the first stage to get predicted prices ($\hat{p}$) using Yelp rating and marginal cost as instruments.
Then, we use these predicted prices in the second stage regression to estimate market share.

From the results in Table \ref{tab:methods_combined}, we find that $\hat{\alpha} = -1.355$, which is negative and significant.
This is more in line with what we would expect, as an increase in price should lead to decreased market share.


\subsection*{Part C: IV Estimation with Marginal Cost and Distance}

We further refine our IV estimation by including average distance (ave\_dist) as an additional instrument.
The estimated equations are:
\begin{gather*}
  s = \beta_0 + \beta_1 x + \beta_2 \hat{p} + \varepsilon \\
  \shortintertext{where} \\
  \hat{p} = \gamma_0 + \gamma_1 x + \gamma_2 mc + \gamma_3 ave\_dist + \nu
\end{gather*}

Similarly, we regress the first stage to get predicted prices ($\hat{p}$) using Yelp rating, marginal cost, and average distance as instruments.
Then, we use these predicted prices in the second stage regression to estimate market share.

From the results in Table \ref{tab:methods_combined}, we find that $\hat{\alpha} = -1.340$, which is again negative and significant.
This further supports the conclusion that price negatively affects market share when accounting for endogeneity. Further,
the magnitude is similar, but a bit smaller, suggesting including distance does help to tune the model some.

\section*{Problem 3: Exogenous Choice Set Formation}

\begin{table}[htbp]
\centering
\caption{GMM Estimates under Exogenous Search}
\label{tab:q3_exogenous}
\begin{tabular}{lccc}
\hline\hline
\textbf{Parameter} & \textbf{Estimate} & \textbf{True Value} \\
\hline
Intercept              & 1.583 & 1.00 \\
Rating Coefficient     & 0.243 & 0.20 \\
Price Coefficient      & -0.543 & -0.25 \\
Mu                     & -0.162 & --- \\
\hline\hline
\end{tabular}
\end{table}

In the exogenous choice set formation model, we assume that consumers consider all available restaurants when making their choices.
Using GMM estimation, we can recover estimates of the parameters governing consumer preferences and choice behavior.

From the results in Table \ref{tab:q3_exogenous}, we observe that the estimated coefficients for rating and price are somewhat close to their true values.
The rating coefficient is estimated at 0.243, which is reasonably close to the true value of 0.20. The price coefficient is estimated at -0.543, which is negative
and significant, aligning with expectations. However, the magnitude is more than double the true value of -0.25, indicating some bias in the estimate. Additionally,
the intercept is estimated at 1.583, which is higher than the true value of 1.00.

This model does not perform well in terms of accurately recovering the true parameters, suggesting that the assumption of exogenous choice set formation
may not be appropriate in this context.

\section*{Problem 4: Endogenous Choice Set Formation}

\begin{table}[htbp]
\centering
\caption{GMM Estimates with Endogenous Search}
\label{tab:q4_endogenous}
\begin{tabular}{lccc}
\hline\hline
\textbf{Parameter} & \textbf{Estimate} & \textbf{True Value} \\
\hline
Intercept              & 2.608 & 1.00 \\
Rating Coefficient     & 0.250 & 0.20 \\
Price Coefficient      & -0.587 & -0.25 \\
Gamma                  & -0.238 & -0.10 \\
\hline\hline
\end{tabular}
\end{table}

In the endogenous choice set formation model, we allow for the possibility that consumers do not consider all restaurants equally.
Instead, choice sets are formed based on certain characteristics, such as distance or Yelp ratings. Using similar methods to those
in Problem 3, we estimate the parameters of this model.

From the results in Table \ref{tab:q4_endogenous}, we observe that the estimated coefficients for rating and price are much closer to their true values.
The rating coefficient is estimated at 0.250, which is reasonably close to the true value of 0.20. The price coefficient is estimated at -0.587, which is negative
and significant, aligning with expectations. The parameter $\hat{\gamma}$, which captures the disutility of search on choice set formation, is also negative,
indicating that greater distance reduces the likelihood of a restaurant being a part of a consumer's choice set.

This model appears to perform much better than the model in problem 3, which leads me to believe that endogenous choice set formation is a much more realistic
assumption about consumer behavior in this context.

\section*{Discussion and Conclusion}

In this problem set, we went through several estimation exercises to better understand consumer choice behavior. We began with basic techniques, OLS and IV, to 
see the issue of endogenity in price and why these linear models may not be sufficient in this context. Then, we moved to more complicated methods, GMM with
exogenous and endogenous choice set formation. The results clearly indicated that accounting for endogenous choice sets provided more accurate estimates of 
consumer preferences. The exogenous case, however, produced questionable results, implying that consumers likely will not take literally every option into account.
Overall, this assignment highlighted the importance of model specification and the need to consider realistic consumer behavior when estimating demand.

\end{document}
